% Insert as A.55.  This is an exact finite correction calculation.
\begin{statement}{A.55}{Proved}{The dual obstruction to correction using completely split auxiliary primes}
Let $N\ge2$, let $M_N=\Z/N(2)$, and identify its Cartier dual by
the evaluation pairing
\[
 M_N^D=\operatorname{Hom}(M_N,\mu_N)=\Z/N(-1),\qquad
 M_N\otimes M_N^D\longrightarrow\mu_N.
 \tag{A.55.1}
\]
Let $K/\Q$ be finite Galois containing $\mu_N$, let $S$ contain
the infinite place, the primes dividing $N$, and the primes ramified
in $K$, and let $T$ be the set of finite primes outside $S$ which
split completely in $K$.  Write $H^1_{S\cup T}(\Q,M_N)$ for the
union of $H^1(G_{\Q,S\cup T_0},M_N)$ over finite $T_0\subset T$.
Then
\[
 \begin{aligned}
 &\operatorname{im}\left(H^1_{S\cup T}(\Q,M_N)
       \longrightarrow\prod_{v\in S}H^1(\Q_v,M_N)\right)
 \\
 &\qquad=\left(\operatorname{res}_S
     H^1(\operatorname{Gal}(K/\Q),M_N^D)\right)^\perp.
 \end{aligned}
 \tag{A.55.2}
\]
Here the annihilator is taken under the actual pairing
\[
 ((z_v),\xi)\longmapsto
       \sum_{v\in S}\operatorname{inv}_v
                        (z_v\cup\operatorname{res}_v\xi)
          \quad\hbox{in }\Q/\Z.
 \tag{A.55.3}
\]

There is a concrete order-sixteen character in this dual group.
Suppose $N=2^r$ with $r\ge4$, take $c=2v^{24}\in\Q^\times$,
and use compatible Kummer coordinates
$b=\kappa_N(c)\in Z^1(G_\Q,\Z/N(1))$.
For $g\in G_\Q$ write
\[
 \chi(g)\bmod16=(-1)^{e(g)}5^{L(g)},\quad L(g)\in\Z/4,
 \qquad
 \beta(g)=\frac{\chi(g)^2-1}{8}\bmod2.
 \tag{A.55.4}
\]
The function $\binom{L}{2}\bmod2$ is well defined on $\Z/4$.
Then
\[
 \zeta_N(g)=\frac N{16}b(g)+\frac N2\binom{L(g)}2
                  \quad\hbox{in }\Z/N(-1)
 \tag{A.55.5}
\]
is a continuous one-cocycle.  Its cohomology class has exact
order $16$ and factors through
\[
                     K_c=\Q(\mu_N,c^{1/N}).
 \tag{A.55.6}
\]
Moreover,
\[
          2\zeta_N=\frac N8\kappa_N(c),\qquad
          8\zeta_N=\frac N2\beta
                  \quad\hbox{as actual cocycles}.
 \tag{A.55.7}
\]
The first expression on the right is viewed in weight $-1$ by
the equivariant coefficient map $\Z/N(1)\to\Z/N(-1)$,
$u\mapsto(N/8)u$.

For an explicit failed correction problem, take
\[
              c=2\cdot73^{24},\qquad S=\{2,73,\infty\},
              \qquad K=K_c.
 \tag{A.55.8}
\]
At $73$ let $z_{73}$ be the unramified class in
$H^1(\Q_{73},M_N)$ with arithmetic Frobenius value $1$,
and prescribe $z_2=z_\infty=0$.  This tuple is not in the
image in (A.55.2), for every $r\ge4$.

There is also an exact constraint on retaining the integral
class-two section.  For every prime $\ell$ and $a\ge1$, the
coefficient connecting map gives the exact sequence
\[
 0\longrightarrow
 \frac{H^1_{\rm cts}(G_\Q,\Z_\ell(2))}{\ell^a}
 \longrightarrow H^1(G_\Q,\Z/\ell^a(2))
 \xrightarrow{\ \partial\ }
 H^2_{\rm cts}(G_\Q,\Z_\ell(2))[\ell^a]
 \longrightarrow0.
 \tag{A.55.9}
\]
A finite central correction $z$ of a fixed integral class-two
lift is the reduction of an integral central correction exactly
when $\partial z=0$.  This is equivalent to $z$ being unramified
outside $\ell$, with the additional condition
$\operatorname{res}_\infty z=0$ when $\ell=2$.
Thus the allowable integral corrections are precisely the
reductions of the groups computed in A.40.  Once its two
cross-prime normalizations are imposed, none of their nonzero
classes survives at modulus divisible by $24$.
\end{statement}
\begin{proof}
\emph{The exact image of correction with split auxiliaries.}
For a finite $T_0\subset T$, apply finite-module Poitou--Tate
duality over $G_{\Q,S\cup T_0}$, in the form
\cite[I, Theorem 4.10(b)]{BIR_MilneADT}.
The annihilator of the image on the $S$ coordinates is
\[
 \operatorname{res}_S\{\xi\in H^1(G_{\Q,S\cup T_0},M_N^D):
                           \operatorname{res}_{T_0}\xi=0\}.
 \tag{A.55.10}
\]
Indeed one extends an $S$-tuple in the dual local groups by
zero at $T_0$ and uses the orthogonality of the two global
images.  A class in braces is unramified at $T_0$, so it
inflates from $H^1(G_{\Q,S},M_N^D)$.  This last group is
finite by \cite[I, Corollary 4.15]{BIR_MilneADT}.
Consequently the decreasing kernels obtained as $T_0$ grows
stabilize, and the annihilator of the union of correction
images is
\[
 \operatorname{res}_S\{\xi\in H^1(G_{\Q,S},M_N^D):
                            \operatorname{res}_q\xi=0\ (q\in T)\}.
 \tag{A.55.11}
\]

We identify the group in braces, including the inflation map.
Since $K$ contains $\mu_N$, $G_K$ acts trivially on $M_N^D$.
The restriction of a representing cocycle to $G_K$ is
therefore a homomorphism.  It factors through a finite Galois
extension of $K$ which can be chosen Galois over $\Q$.
Chebotarev, applied to every conjugacy class in the kernel
over $K$, realizes that class by a prime in $T$.  At such a
prime the module is trivial and the unramified local class
is its Frobenius value.  Local vanishing thus makes the
homomorphism on $G_K$ zero.  Inflation--restriction now gives
exactly the image of
$H^1(\operatorname{Gal}(K/\Q),M_N^D)$.
Conversely every inflated class vanishes on each such local
decomposition group and is unramified outside $S$.
The finite perfect local pairings identify a subgroup with
its double annihilator, proving (A.55.2).  Formula (A.55.3)
is the cup-product map in this argument, followed by the
local Brauer invariant; no comparison of different coefficient
modules is implicit.

\emph{The order-sixteen cocycle.}
The identity $\sqrt2=\zeta_8+\zeta_8^{-1}$ gives
\[
                   b\bmod2=\kappa_2(2)=\beta=L\bmod2.
 \tag{A.55.12}
\]
This is an equality of cocycles: the coefficient action
modulo two is trivial and hence its one-coboundaries vanish.
The map $L:G_\Q\to\Z/4$ is a homomorphism.  Its binomial
cochain satisfies
\[
 d\binom L2(g,h)=-\beta(g)\beta(h)\quad\hbox{modulo }2.
 \tag{A.55.13}
\]
Write $d_{-1}$ for the differential in weight $-1$.
Since $b$ is a cocycle in weight $1$,
\[
 d_{-1}\left(\frac N{16}b\right)(g,h)
 =\frac N{16}(\chi(g)^{-1}-\chi(g))b(h)
 =-\frac N2\beta(g)\beta(h)\quad\hbox{modulo }N.
 \tag{A.55.14}
\]
The differential of the second summand of (A.55.5) is the
same order-two value.  Their sum is zero, proving that
$\zeta_N$ is a cocycle.  All its data factor through (A.55.6).
Multiplication by two and by eight gives (A.55.7).

The cochain $16\zeta_N$ is zero.  To prove that
$8\zeta_N$ is not a coboundary, suppose it were
$(\chi^{-1}-1)m$ for some $m\in\Z/N$.
At complex conjugation, its value is zero, so $2m=0$.
Thus $m=0$ or $N/2$.  At an element with
$\chi\equiv5\pmod {16}$ the coboundary of either choice
is zero, while $(N/2)\beta=N/2$.  This contradiction proves
exact order sixteen.  The equivariance of the map used in
(A.55.7) follows from $8\mid u^2-1$ for every odd $u$.

\emph{The local pairing giving the counterexample.}
The extension in (A.55.8) is unramified outside $S$.
At $q=73$, the tame inertia value of $b$ is $24$, and the
cyclotomic binomial cochain is unramified.  Hence
\[
                        \zeta_N(T_{73})=N/2.
 \tag{A.55.15}
\]
An unramified weight-two class with Frobenius value $1$
exists because
$H^1_{\rm ur}(\Q_q,M_N)=M_N/(q^2-1)M_N$.
For an unramified weight-two cocycle $z$ of Frobenius
value $d$, paired with a weight-minus-one cocycle of inertia
value $t$, the tame residue of their cup is $qdt$ in
$\Z/N$.  This follows directly from the strict formula
of A.27: the $(\Phi,T)$ cup term is $dq^2t$, its two
middle terms cancel, its inertia terms vanish, and the
residue map multiplies by $q^{-1}$.
In (A.55.15), $q$ is odd, so the residue is $N/2$.
Consequently
\[
              \operatorname{inv}_{73}(z_{73}\cup\zeta_N)
                                 =\frac12\quad\hbox{in }\Q/\Z.
 \tag{A.55.16}
\]
At the other two members of $S$ the prescribed class is zero.
Thus (A.55.3) is nonzero and (A.55.2) excludes a global
correction with the allowed ramification.  Equivalently,
global Brauer reciprocity alone yields the contradiction:
at $T$ the dual class is zero, and outside $S\cup T$ both
classes are unramified, whose cup has zero invariant.

\emph{Preservation of the integral central lift.}
The continuous cochain sequence for
$0\to\Z_\ell(2)\xrightarrow{\ell^a}\Z_\ell(2)
\to\Z/\ell^a(2)\to0$ gives (A.55.9).
Adding a finite cocycle to the central coordinate $h$ is
the actual action on the class-two lifts of fixed abelian
cocycles $a,b$; reduction commutes with this action.
This proves the stated Bockstein criterion on lifts.

Here is its exact ramification test.  At $q\ne\ell$, put
$t=z(T_q)\in\Z/\ell^a$.  The tame cocycle relation is
$(q-1)t=0$.  Choose an integral representative $\widetilde t$.
The strict cochain identity of A.27 sends the local Bockstein
to
\[
       \frac{(q-1)\widetilde t}{\ell^a}
                    \pmod{(q-1)\Z_\ell}
       \quad\hbox{in }H^2(\Q_q,\Z_\ell(2)).
 \tag{A.55.17}
\]
This class vanishes exactly when $t=0$.  Thus every
integrally liftable correction is unramified outside $\ell$.
Conversely, such an unramified cocycle descends to
$G_{\Q,\{\ell,\infty\}}$: it kills the normal closure of
the other inertia groups, on which the coefficients are
trivial.  For odd $\ell$, A.27 gives vanishing of the
unramified integral $H^2$, and hence vanishing of its
Bockstein.  At two that $H^2$ is $\Z/2$ and restriction
to the real place is injective.  At the real place
$H^1(G_\R,\Z_2(2))=0$, so the finite-class Bockstein
vanishes exactly when the finite local class is zero.
This proves both directions of the test.  A.40 identifies
its global kernel and its cross-prime restrictions.

The finite certificate
\texttt{split\_auxiliary\_certificate.py} checks the
twisted cocycle identity exhaustively modulo sixteen,
the order and coefficient comparisons, the local pairing,
and the integral Bockstein residue test.  Scaling the
modulo-sixteen identity by $N/16$ proves the same identity
at every $N=2^r$, $r\ge4$.
\end{proof}

For the actual ghosts of A.44, a finite representative of
$-B$ has the form $2v^{24}$.  The same formula therefore
supplies a member of the dual correction group for any
auxiliary field containing its Kummer field and $\mu_N$.
At a support prime $r_j$ with $v_2(n_j)=3$, its inertia
value is $N/2$, so it tests the parity of an unramified
local central correction there.  Formula (A.55.2) gives
the exact test for any proposed tuple of corrections.
The tuple specified in (A.55.8) is an explicit failed
correction problem; the calculations above do not identify
it with the as-yet uncomputed local discrepancy tuple of
the distinguished ghost lift $s_*$.  That tuple must still
be evaluated using A.40 and A.46.
